\chapter{Mar.~4 --- Affine Lie Algebras}

\section{Affine Lie Algebras, Continued}
\begin{remark}
  We can extend
  $\widehat{\g}$ again
  by the grading operator
  $d = t \frac{d}{dt}$ to get
  $\widetilde{\g} = \widehat{\g} \oplus \C d$,
  where
  \[
    [d, x[n]] = n x[n] \quad \text{and} \quad
    [d, x \otimes P] = x \otimes t\frac{t}{dt} P.
  \]
  We can also define a nondegenerate form
  on $\widehat{g}$ by
  \[
    (x \otimes P, y \otimes Q)
    = \frac{(x, y)}{2\pi i} \oint_{|t| = 1} P Q\, \frac{dt}{t},
  \]
  and on $\widetilde{\g}$ by
  $(c, d) = 1$ and $(c, c) = (d, d) = (c, x \otimes P) = (d, x \otimes P) = 0$.
\end{remark}

\begin{definition}
  Define $\widehat{\mathfrak{h}} \subseteq \widehat{\g}$
  by $\widehat{\mathfrak{h}} = \mathfrak{h} \oplus \C c$
  and $\widetilde{\mathfrak{h}} \subseteq \widetilde{\g}$
  by $\widetilde{\mathfrak{h}} = \widehat{\mathfrak{h}} \oplus \C d$,
  these are Cartan subalgebras.
  Their duals are
  \[
    \widehat{\mathfrak{h}}^* = \mathfrak{h}^* \oplus \C \Lambda_0
    \quad \text{and} \quad
    \widetilde{\mathfrak{h}}^* = \mathfrak{h}^* \oplus \C \Lambda_0 \oplus \C \delta,
  \]
  where $\Lambda_0(c) = 1$,
  $\Lambda_0(d) = \Lambda_0(h) = 0$ for
  $h \in \mathfrak{h}$,
  and $\delta(d) = 1$,
  $\delta(c) = \delta(h) = 0$ for $h \in \mathfrak{h}$.
\end{definition}

\begin{remark}
  We can write
  \[
    \widetilde{\g} = \widetilde{\mathfrak{h}}
    \oplus \bigoplus_{\gamma \in \widehat{R}}
    \widehat{\g}^\gamma,
  \]
  where $\widetilde{\g}^\gamma = \{x \in \widetilde{\g} : [h, x] = \gamma(h) x \text{ for } h \in \widetilde{\mathfrak{h}}\}$
  and $\widehat{R} = \{\alpha + n \delta : \alpha \in R \sqcup \{0\}, n \in \Z \text{ not both $0$}\}$.
  Then
  \[
    \widetilde{\g}^{\alpha + n \delta}
    = \g^\alpha \otimes t^n
    \quad \text{and} \quad
    \widetilde{\g}^{n \delta} = \mathfrak{h} \otimes t^n,
  \]
  with $\alpha \in R$, $n \in \Z$, and
  $n \in \Z \setminus \{0\}$, respectively.
  We can choose a polarization
  \[
    \widehat{R} = \widehat{R}_+ \sqcup \widehat{R}_-,
    \quad \widehat{R}_+ = \{\alpha + n \delta : n > 0 \text{ or } n = 0, \alpha \in R_+\},
  \]
  where as usual $\widehat{R}_- = - \widehat{R}_+$.
  The simple roots in $\widehat{R}_+$ are
  the ones which cannot be represented
  as a nontrivial sum of positive roots.
  Let $\alpha_1, \dots, \alpha_r$ be
  the simple roots of $\g$ and
  \[
    \alpha_0 = \delta - \theta,
  \]
  where $\theta$ is the highest root. Then
  we can define
  $\widehat{\mathfrak{n}}^{\pm}
    = \mathfrak{n}^{\pm}
    \oplus (\g \otimes t^{\pm 1} \C[t^{\pm 1}])$
  and write
  \[
    \widetilde{\g}
    = \widehat{\mathfrak{n}}^+ \oplus \widehat{\mathfrak{h}} \oplus \widehat{\mathfrak{n}}^-.
  \]
  The root lattices $\widehat{Q} = \bigoplus_{i = 0}^r \Z \alpha_i$
  and $\widehat{Q}^+= \sum_{i = 0}^r \Z_{\geq 0} \alpha_i$.
  We can define the coroot
  \[
    \alpha_0^\vee
    = \frac{2h_{\alpha_0}}{(\alpha_0, \alpha_0)}
    = h_{\alpha_0}
    = c - \theta^\vee.
  \]
  Then we can define the weight lattice by
  $\widehat{P} = \langle \Lambda \in \widetilde{\mathfrak{h}}^* : \Lambda(\alpha_i^\vee) \in \Z \rangle = P \oplus \Z \Lambda_0 \oplus \C \delta$, and
  \[
    \widehat{P}^+
    = \{\Lambda \in \widetilde{\mathfrak{h}}^* : \Lambda(\alpha_i^\vee) \in \Z_+\}
    = \{\Lambda = \lambda + k \Lambda_0 - \Delta \delta : \lambda \in P^+, k \in \Z_+, \Delta \in \C, \langle \lambda, \theta^\vee \rangle \le k\}.
  \]
  From this we get a Cartan matrix
  $A = \{a_{i, j}\}_{i, j = 0}^r$.
\end{remark}

\section{Generalized Cartan Matrices}
\begin{remark}
  A \emph{generalized Cartan matrix}
  is a matrix $A = \{a_{i, j}\}_{i, j = 0}^r$
  satisfying
  \begin{enumerate}
    \item $a_{i, i} = 2$ for all $i$,
    \item $a_{i, j} \le 0$ for $i \ne j$,
    \item $a_{i, j} = 0$ if and only if
      $a_{j, i} = 0$.
  \end{enumerate}
\end{remark}

\begin{definition}
  Let $A$ be a generalized Cartan matrix.
  The \emph{Kac-Moody algebra} $\g(A)$
  is the Lie algebra generated by elements
  $\{e_i, h_i, f_i\}$ together with the
  standard Serre relations.
\end{definition}

\begin{theorem}
  The Lie algebra $\widehat{\g}$
  is the Kac-Moody algebra associated with
  the Cartan matrix which we constructed
  above.
\end{theorem}

\begin{remark}
  We have
  $e_0 \in \g^{-\theta} \otimes t$,
  $f_0 \in \g^\theta \otimes t^{-1}$,
  $(e_0, f_0) = 1$, and
  $h_0 = [e_0, f_0] = c - \theta^\vee$.
\end{remark}

\begin{remark}
  Let $A$ be a generalized Cartan matrix and
  assume that it is symmetrizable, i.e.
  there exist $d_i$ such that
  $d_i a_{i, j} = d_j a_{j, i}$. Then
  we can define an inner product
  \[
    (h_i, h_j) := \frac{a_{i, j}}{d_i}
  \]
  on the Cartan subalgebra. Let
  $Z$ be the kernel of this inner product.
  Let
  \[\g_e(A) = \g(A) \oplus Z^*\]
  with commutation relations
  $[z_1, z_2] = [z_1, h] = 0$,
  $[z, e_j] = z(\alpha_i) e_i$,
  and $[z, f_j] = -z(\alpha_j) f_j$. Then
  \[
    \langle h_i, h_j \rangle = \frac{\alpha_{i, j}}{d_j},
    \quad \langle h, z \rangle = z(h), \quad
    \langle z_1, z_2 \rangle = 0
  \]
  defines a nondegenerate pairing
  on $\g_e(A)$.
\end{remark}

\section{Verma Modules and Weyl Modules}

\begin{definition}
  Let $\Lambda \in \widetilde{\mathfrak{h}}^*$.
  The \emph{Verma module}
  corresponding to $\Lambda$ is
  \[
    M_\Lambda
    = \Ind_{\widehat{\mathfrak{h}} \oplus \widehat{\mathfrak{n}}^+}^{\widetilde{\g}} \C_\Lambda
  \]
  where $\C_{\Lambda}$ is the $1$-dimensional
  representation of
  $\widehat{\mathfrak{h}} \oplus \widehat{\mathfrak{n}}^+$.
  Let $v_\Lambda$ be the basis vector, with
  $h v_\Lambda = \Lambda(h) v_\Lambda$.
\end{definition}

\begin{remark}
  We have a weight decomposition given by
  \[
    M_\Lambda = \bigoplus_{m \in \Lambda - \widehat{Q}^+} M_\Lambda^m,
  \]
  where $\Lambda = \lambda + k \Lambda_0 - \Delta \delta$
  for some $\lambda \in \mathfrak{h}^*$.
  From the $\widehat{g}$ perspective, we
  have
  \[
    \Delta = \Delta(\lambda)
    = \frac{(\lambda, \lambda + 2 \rho)}{2(k + h^\vee)}
  \]
  (here $h^\vee$ is the dual Coxeter number).
  We call $k$ the \emph{central charge}
  (or \emph{level}) of the representation.
  In this case, we say that a
  $\widehat{g}$-module $V$ is
  of level $k$ or $V$ is a module with
  central charge $k$. We denote the
  corresponding Verma module $M_{\Lambda}$
  by $M_{\lambda, k}$.

  The case
  $k = - h^\vee$ is called the
  \emph{critical level}, though this
  generally does not happen.
\end{remark}

\begin{remark}
  The Verma module $M_{\lambda, k}$ has
  a natural $\Z$-grading given by
  \[
    M_{\lambda, k}
    = \bigoplus_{n \ge 0} M_{\lambda, k}[-n],
  \]
  where $M_{\lambda, k}[-n]$ is naturally a
  $\g$-module and an eigenspace of
  $d$ with eigenvalue $-n - \Delta(\lambda)$.
\end{remark}

\begin{example}
  We have $M_{\lambda, k}[0] = M_\lambda$,
  the Verma module for $\g$.
\end{example}

\begin{remark}
  Let $L_\Lambda = M_\Lambda / I_\Lambda$,
  where $I_\Lambda$ is the maximal
  proper submodule of $M_\Lambda$. Then
  $L_{\lambda, k} := L_\Lambda$ is an irreducible module.
  The \emph{Chevalley involution}
  $\widetilde{\omega}$
  is given by
  \[
    \widetilde{\omega}(d) = -d
  \]
  and on $L\g = \g \otimes \C[t, t^{-1}]$
  by $t \mapsto t^{-1}$. Then
  $\widetilde{\omega}$ gives a
  contravariant \emph{Shapovalov form}
  with kernel $I_\Lambda$.
\end{remark}

\begin{remark}
  We can construct a class of
  representations of $\widetilde{\g}$ induced
  from $\g$-modules. Let
  \[
    \widetilde{g}^+
    = \g \otimes \C[t] \oplus \C c \oplus \C d \subseteq \widetilde{\g}.
  \]
  If $V$ is a module for $\g$,
  then automatically $V$ is also a module
  for $\widetilde{\g}^+$ by letting
  \begin{enumerate}
    \item $\g \otimes t\C[t]$ act by $0$,
    \item $c, d$ act by
      $k, -\Delta$, respectively.
  \end{enumerate}
  Then one can define
  $\Ind_{\widetilde{g}^+}^{\widetilde{\g}} V$,
  which is a module for $\widetilde{\g}$.
\end{remark}

\begin{example}
  Verma modules belong to this class,
  as we can write
  \[
    M_{\lambda, k}
    = \Ind_{\widetilde{\g}^+}^{\widetilde{\g}} M_\lambda
  \]
  where $d$ acts on $M_\lambda$ by
  $-\Delta(\lambda)$.
\end{example}

\begin{example}
  Let $V_{\lambda, k} = \Ind_{\widetilde{\g}^+}^{\widetilde{\g}} L_\lambda$,
  where $d$ acts as $-\Delta(\lambda)$.
  If $\lambda \in P_+$, we call $V_{\lambda, k}$
  a \emph{Weyl module}.
\end{example}

\begin{remark}
  We get a contragradient Verma module by
  replacing $\omega$ with $\widetilde{\omega}$.
  This gives maps
  $M_{\lambda, k} \to M_{\lambda, k}^c$ and
  $V_{\lambda, k} \to V_{\lambda, k}^c$,
  which are isomorphisms if
  $M_{\lambda, k}, V_{\lambda, k}$ are
  irreducible.
\end{remark}

\section{Integrable Representations}

\begin{theorem}
  The Verma module $M_\Lambda$ is
  irreducible if the following
  condition is satisfied:
  For every $\widehat{\alpha} \in \widehat{R}_+$
  and $N \in \{1, 2, \dots\}$, we have
  $(\Lambda + \widehat{\rho}, \widehat{\alpha}) \ne N(\widehat{\alpha}, \widehat{\alpha}) / 2$,
  where $\widehat{\rho} = \rho + h^\vee \Lambda_0$.
\end{theorem}

\begin{remark}
  For $(\widehat{\alpha}, \widehat{\alpha}) \ne 0$,
  the above condition becomes
  $(\Lambda + \widehat{\rho}, \widehat{\alpha}^\vee) \ne N$,
  which is the analogue of the
  simple Lie algebra case.
\end{remark}

\begin{remark}
  We call weights satisfying
  $(\Lambda + \widehat{\rho}, \widehat{\alpha}) \ne N(\widehat{\alpha}, \widehat{\alpha}) / 2$
  \emph{generic}.
  For generic weights, $L_\Lambda = M_\Lambda$.
\end{remark}

\begin{theorem}
  Let $\lambda \in \mathfrak{h}^*$ and
  $k \in \C$ be such that
  $k \notin \Q \langle \alpha^\vee, \lambda \rangle + \Q$
  for all $\alpha \in R_+$. Then the
  induced module $V_{\lambda, k}$
  is irreducible.
  If this holds, we say that
  $k$ is \emph{generic with respect to $\lambda$}.
\end{theorem}

\begin{corollary}
  If $\lambda \in P^+$ and
  $k \notin \Q$, then the Weyl module
  $V_{\lambda, k}$
  is irreducible.
\end{corollary}
